% !TeX root = ../papers/nc-p3.tex
\section{第二長方形の階数十四・周期二曲線の全点分類}
\label{sec:rank-fourteen-conic-complete}
\begin{theorem}\label{thm:rank-fourteen-conic-complete}
命題\ref{prop:second-rectangle-rank-fourteen-period}の対角曲線の全点は標準Koszul左右AS正則で、
最小周期2、射影安定群は四元群である。指定始点の無印同型類は射影直線から十二点を除いた曲線と一対一に対応する。
座標を $s=h^2$ とすれば、除くのは $0,\infty$ および
\[
 (s^2+8s+4)(s^2-8s+4)(s^2+4)
 \big((s^2+8s+4)^2-16s(s-2)^2\big)=0
\]
の根である。残る点のうち十二点に一つの厚い係数の消滅があり、他では台147の全52係数が非零である。
一位置移動は $s\mapsto-s$ である。
\end{theorem}

\begin{proof}[パラメータと枠]
$i^2=-1$ とし、前節の円錐曲線で $h=S+iE$ と置く。
$S=(h-2/h)/2$, $E=(h+2/h)/(2i)$ である。
$A=h^2+2h-2$, $B=h^2-2h-2$, $C=h^2+2ih+2$, $D=h^2-2ih+2$ とすると
\[
 u=-AC/(BD),\qquad v=-AD/(BC),\qquad z=w=i(h^4-4)/(4h^2).
\]
$r=0,1,-1$ で逆変換が特異となる可能性については、有限な $l,m$ で $r=0$ は
円錐曲線の元の方程式に反し、$r=\pm1$ なら $l^2=m^2=1$ となり $u,v$ はトーラスにない。
$u=-1$ または $v=-1$ は $D(A',B')=0$（$A'=u+1/u,B'=v+1/v$）において他方も $\pm1$ を強制し、
元の階数14層には入らない。$h=0,\infty$ の極限も $u=v=-1$ である。
トーラス開条件の分子・分母の平方自由積は
\[
 h(h^4+8h^2+4)\prod_{e=\pm1}(h^2+2eh-2)(h^2+2eh+2)(h^4+4eh^3+8h^2-8eh+4).
\]
従って有限遠で21点、無限遠を加えて22点が除かれる。

テンソルの交代勾配を使い、各位置で他の三位置に対する三つの二次モーメントの対角成分を作る。
共通分母を払った後、任意の二文字の三つの値の差について $\mathbb Q(i)[h]$ から
$\mathbb Q[h]$ へのノルムの最大公約数を取る。上の境界の全因子を取り除くと四位置の全24対で定数になる。
よって共同固有線は全点で四つの元の文字線であり、印を忘れた同型は単項変換である。

開曲線上で消え得る厚い係数は $x_3,x_6,x_8,x_{10}$ だけである。各々で
境界を除いた零因子は $\mathbb Q(i)$ 上の相異なる二つの三次既約式で、計八式は全て互いに素である。
このため一点で消えるのは一係数だけであり、指定した係数枠で24点が該当する。
これら一欠損台の間の四置換を全探索すると、全て元の52成分の台の置換に延びる。
薄い係数のゲージと $B=1$ の規格化は引き続き可逆なので、
定理\ref{thm:parallel-horizontal-moduli}の $I,J$ がそのまま全点で使える。
対角倍率の非零Smith因子は各台で $1^{11},2,2$ である。
$I$ は $h\mapsto-h$、一位置移動 $J$ は $h\mapsto ih$ となる。
$h\ne0,\infty$ より指定始点で非自明な台置換は点を固定せず、安定群は対角四元群である。
よって指定始点の商は $h^2$ であり、最小周期は2である。
\path{computations/verify_rank_fourteen_conic_frames.py} が境界、係数の零点、共同固有線と全台置換を検算する。
\end{proof}

\begin{proof}[全点の正則性]
始点0で五中心商を計算し、定理\ref{thm:first-rectangle-edge-regularity}の五変数上の表示を使う。
表示Aは位置別文字移動 $(0,0,2,2)$ と第一中心座標の定数1で、全52係数の非零を仮定する。
表示Bは $(2,2,0,0)$ と第二中心座標の定数1で、境界22点だけを除いた基底環を使う。
係数を $(a+ib)/d$, $a,b,d\in\mathbb Q[h]$ で表し、ピボットの分子ノルム $a^2+b^2$ と
分母の既約因子を記録する。前節の単元小行列で特殊化する関係フレームを取り、二次・四次の経路関係を順に消去する。
両表示で全中心商の長さ別基底個数は $(1,4,10,20,30,36,34,20,0)$、
必要な独立関係は次数9から16まで $20,34,36,30,20,10,4,1$ 個である。
評価後の関係の階数条件までに使う境界外の既約因子は各表示で174個と171個である。
最初の表示に追加で仮定した全係数非零の条件も含め、同時に両表示から漏れるのは
次数 $2,2,4,4,6,6,6,6,6,6,8,12,16,16$ の十四個の有理既約式の零点に限られる。

各式について次数1の単純剰余根を持つ素数 $p\equiv1\pmod4$ を選び、
$i$ の二つの剰余を\emph{ともに}使って四中心商を長さ13まで全始点で検算する。
有理既約式の一根で定める数体と $\mathbb Q(i)$ の合成は、直積になる場合でもこの二つの選択で
全ての体因子が覆われる。元の経路の全次数下界と持ち上げた中心商の上界から四中心式上の自由性が従う。
曲線は $J$ で保たれるので、開集合の始点0の結果も全点の全始点に及ぶ。
\path{computations/verify_rank_fourteen_conic_presentation.py} と
\path{research/rank_fourteen_conic_residues.json} に両表示と十四の残余検証を与えた。
\end{proof}
